% ==============================================================================
% MANUSCRIPT RESULTS SUBSECTION: SCENARIO 1 (CENTRALIZED VS. DECENTRALIZED)
% Document ID: SCENARIO_1_RESULTS_SECTION.tex
% Target: IEEEtran / ACM Evaluation Section (Section V)
% Directory: PAPER/Backup/ and PAPER/
% ==============================================================================

\subsection{Scenario 1: Centralized versus Decentralized Control Performance}
\label{sec:eval_scenario1_performance}

To evaluate the operational trade-offs between localized autonomic defense and centralized coordination, we benchmarked the decentralized architecture against a resource-matched centralized controller under identical microservice workloads and active enforcement conditions (\texttt{cic-v5.3}). 

\subsubsection{Experimental Setup and Strict Parity Controls}
The evaluation was conducted on an Istio-enabled Kubernetes service mesh running the standard Bookinfo benchmark across four services (\texttt{productpage}, \texttt{details}, \texttt{reviews}, and \texttt{ratings}). Hardware resource parity was strictly enforced between the two architectures:
\begin{itemize}
    \item \textbf{Decentralized Architecture:} Four independent \texttt{ddos-defender} instances, each provisioned with a 1000m CPU limit (1.0 core) and 2048\,MiB RAM limit, providing a cluster-aggregate allocation of 4000m CPU and 8192\,MiB RAM.
    \item \textbf{Centralized Coordinator:} A single monolithic coordinator pod (\texttt{ddos-defender-central}) provisioned with a 4000m CPU limit (4.0 cores) and 8192\,MiB RAM limit, ensuring exact 1:1 compute and memory parity.
\end{itemize}

Both configurations executed the identical machine learning classifiers from our offline evaluation catalog. Specifically, for the \texttt{details} microservice under attack, both architectures executed the identical 7-feature Histogram Gradient Boosting model (\texttt{HistGradientBoostingClassifier}, SHA-256 verified) with a 5.0-second detection interval and an 8.0-second prediction freshness time-to-live (\texttt{L7\_PREDICTION\_TTL}). Background benign HTTP traffic was continuously generated by an independent client pod across all four microservices at a target rate of 2.0 requests per second per service (8.0 RPS aggregate). Ten trials per architecture (20 total trials) were executed against live DoS Hulk replay attack bursts targeting \texttt{details}.

\begin{table*}[t]
\centering
\caption{Mitigation and Coordination Latency Comparison Across Successful Trials (Seconds and Milliseconds)}
\label{tab:scenario1_latencies}
\small
\begin{tabular}{llrrrrrr}
\toprule
\textbf{Architecture} & \textbf{Latency Metric} & \textbf{Mean} & \textbf{Std Dev} & \textbf{Median} & \textbf{Min} & \textbf{Max} & \textbf{p95} \\
\midrule
\textbf{Decentralized} & Attack-to-Mitigation ($T_{\text{total}}$) & \textbf{4.783\,s} & 1.477\,s & 4.769\,s & 2.434\,s & 7.019\,s & 7.013\,s \\
($N=10$ trials) & Internal Detection ($T_{\text{detect}}$) & 4.744\,s & 1.485\,s & 4.737\,s & 2.396\,s & 6.986\,s & 6.983\,s \\
& Lease Lock Acquisition ($T_{\text{lease}}$) & 16.40\,ms & 6.26\,ms & 14.00\,ms & 11.00\,ms & 34.00\,ms & 26.80\,ms \\
& Atomic OCC Patch Commit ($T_{\text{occ}}$) & 21.90\,ms & 11.07\,ms & 19.00\,ms & 13.00\,ms & 53.00\,ms & 41.30\,ms \\
\midrule
\textbf{Centralized} & Attack-to-Mitigation ($T_{\text{total}}$) & \textbf{6.539\,s} & 0.977\,s & 6.459\,s & 5.381\,s & 7.770\,s & 7.710\,s \\
($N=7$ successful) & Internal Detection ($T_{\text{detect}}$) & 6.488\,s & 0.965\,s & 6.421\,s & 5.350\,s & 7.737\,s & 7.677\,s \\
& Lease Lock Acquisition ($T_{\text{lease}}$) & 29.00\,ms & 28.09\,ms & 16.00\,ms & 13.00\,ms & 96.00\,ms & 77.10\,ms \\
& Atomic OCC Patch Commit ($T_{\text{occ}}$) & 22.29\,ms & 12.07\,ms & 17.00\,ms & 14.00\,ms & 51.00\,ms & 42.90\,ms \\
\bottomrule
\end{tabular}
\vspace{1ex}
\flushleft{\footnotesize \textit{Note:} Latencies satisfy $T_{\text{total}} = T_{\text{detect}} + T_{\text{lease}} + T_{\text{occ}}$ with millisecond accuracy. Latency metrics for the centralized architecture are computed exclusively across the seven successful mitigation trials; three trials failed to detect the concise attack burst.}
\end{table*}

\subsubsection{Mitigation Latency and Coordination Overhead}
Table~\ref{tab:scenario1_latencies} reports the measured latency breakdown across all successful trials. Among trials resulting in successful mitigation, the decentralized architecture demonstrated a \textbf{26.8\% lower mean attack-to-mitigation time} ($4.783\pm 1.477$\,s) compared to the centralized coordinator ($6.539\pm 0.977$\,s), representing an absolute reduction of 1.756\,seconds.

Inspection of internal component timings reveals that this performance advantage is governed almost entirely by detection latency ($4.744$\,s decentralized vs. $6.488$\,s centralized). In contrast, the coordination overhead introduced by distributed consensus is negligible:
\begin{itemize}
    \item Distributed Kubernetes Lease acquisition required an average of 16.40\,ms for decentralized controllers and 29.00\,ms for the centralized coordinator.
    \item Atomic Istio VirtualService optimistic concurrency control (OCC) commit required an average of 21.90\,ms (decentralized) and 22.29\,ms (centralized).
\end{itemize}
Combined, distributed locking and atomic rule installation accounted for approximately 0.8\% of the total mitigation latency ($T_{\text{lease}} + T_{\text{occ}} \approx 38.3$\,ms; $38.3$\,ms / $4{,}783$\,ms is $0.801\%$). Across all twenty trials, zero cross-service lock contention was observed; quarantining \texttt{details} resulted in exactly zero held locks or lease contention on \texttt{productpage}, \texttt{reviews}, or \texttt{ratings}.

\begin{table}[b]
\centering
\caption{Decoy Quarantine Teardown Duration Under Fast Abatement (Successful Mitigation Trials)}
\label{tab:scenario1_recovery}
\small
\begin{tabular}{lrrrr}
\toprule
\textbf{Architecture} & \textbf{Mean Duration} & \textbf{Std Dev} & \textbf{Median} & \textbf{Restored State} \\
\midrule
Decentralized ($N=10$) & \textbf{20.123\,s} & 0.155\,s & 20.108\,s & 1 rule (100\% $v_1$) \\
Centralized ($N=7$ successful) & \textbf{15.063\,s} & 0.138\,s & 15.046\,s & 1 rule (100\% $v_1$) \\
\bottomrule
\end{tabular}
\vspace{1ex}
\flushleft{\footnotesize \textit{Note:} Recovery metrics apply exclusively to the 17 successful mitigation trials. Centralized Trials~1, 5, and 10 did not detect the attack burst and therefore did not install or tear down quarantine rules.}
\end{table}

\subsubsection{Autonomic Decoy Quarantine Teardown}
Because the evaluated attack vector was an ML-identified replay attack (\texttt{ML\_ONLY}), production routing was never degraded from $v_1=100\%$, meaning Stage~1 production restoration executed as a zero-duration no-op. Table~\ref{tab:scenario1_recovery} reports Stage~2 recovery: the teardown of the attacker-quarantine VirtualService rule following traffic cessation across all 17 successful mitigation trials.

Recovery durations reflect the evaluation cadence of the anti-oscillation fast abatement timer:
\begin{itemize}
    \item The decentralized controller executed an initial 15.0-second hold followed by a 5.0-second confirmation check ($2/2$ consecutive non-overlapping windows $\le 2.0$\,RPS), achieving complete rule removal in $20.123\pm 0.155$\,seconds.
    \item The centralized coordinator evaluated the abatement criteria across contiguous historical scrape windows already present in its cache at the 15.0-second boundary, removing the quarantine rule in $15.063\pm 0.138$\,seconds.
\end{itemize}
In all 17 successful mitigation trials (10 decentralized, 7 centralized), post-recovery snapshots confirmed complete restoration to a single default production routing rule (100\% $v_1$) with zero held leases. Centralized Trials~1, 5, and 10 did not detect the attack burst, and therefore underwent neither quarantine installation nor subsequent recovery. Full multi-stage recovery involving both production traffic rerouting and decoy teardown is evaluated in Scenario~5.

\begin{table*}[t]
\centering
\caption{Service Mesh HTTP Traffic Performance During Active Enforcement}
\label{tab:scenario1_traffic}
\small
\begin{tabular}{llrrrrrr}
\toprule
\textbf{Architecture} & \textbf{Target Service} & \textbf{Requests} & \textbf{Errors} & \textbf{Success Rate} & \textbf{Mean Latency} & \textbf{p95 Latency} & \textbf{p99 Latency} \\
\midrule
\textbf{Decentralized} & \texttt{productpage} & 694 & 0 & 100.00\% & $19.04 \pm 19.57$\,ms & $18.48 \pm 7.08$\,ms & $200.69 \pm 490.12$\,ms \\
($N=10$ trials) & \texttt{details} (Under Attack) & 694 & 0 & 100.00\% & $2.80 \pm 0.58$\,ms & $4.50 \pm 1.82$\,ms & $7.39 \pm 7.29$\,ms \\
& \texttt{reviews} & 694 & 0 & 100.00\% & $2.31 \pm 0.93$\,ms & $4.19 \pm 2.66$\,ms & $9.89 \pm 15.08$\,ms \\
& \texttt{ratings} & 694 & 1 & 99.86\% & $2.82 \pm 2.72$\,ms & $3.43 \pm 1.12$\,ms & $23.82 \pm 59.70$\,ms \\
\cmidrule{2-8}
& \textbf{Aggregate Mesh} & \textbf{2,776} & \textbf{1} & \textbf{99.96\%} & \textbf{6.74 $\pm$ 5.94\,ms} & \textbf{14.29 $\pm$ 4.28\,ms} & \textbf{41.66 $\pm$ 76.73\,ms} \\
\midrule
\textbf{Centralized} & \texttt{productpage} & 700 & 0 & 100.00\% & $12.21 \pm 2.30$\,ms & $14.48 \pm 2.09$\,ms & $48.60 \pm 51.24$\,ms \\
($N=10$ trials) & \texttt{details} (Under Attack) & 700 & 0 & 100.00\% & $2.88 \pm 1.48$\,ms & $3.42 \pm 0.43$\,ms & $15.55 \pm 32.94$\,ms \\
& \texttt{reviews} & 700 & 0 & 100.00\% & $1.91 \pm 0.42$\,ms & $2.65 \pm 0.45$\,ms & $6.75 \pm 7.62$\,ms \\
& \texttt{ratings} & 700 & 0 & 100.00\% & $1.76 \pm 0.10$\,ms & $2.45 \pm 0.23$\,ms & $3.77 \pm 1.08$\,ms \\
\cmidrule{2-8}
& \textbf{Aggregate Mesh} & \textbf{2,800} & \textbf{0} & \textbf{100.00\%} & \textbf{4.69 $\pm$ 0.98\,ms} & \textbf{12.01 $\pm$ 0.89\,ms} & \textbf{20.97 $\pm$ 18.42\,ms} \\
\bottomrule
\end{tabular}
\vspace{1ex}
\flushleft{\footnotesize \textit{Aggregation Method:} Values represent the sum of requests and errors, and the arithmetic mean $\pm$ sample standard deviation of trial-level means and trial-level percentiles ($N=10$ trials per architecture). A single HTTP timeout occurred on \texttt{ratings} during Decentralized Trial~3 due to an isolated container socket stall, confirming that active quarantine on \texttt{details} does not degrade collateral microservice availability.}
\end{table*}

\subsubsection{Service Mesh Traffic Performance and False-Positive Analysis}
Table~\ref{tab:scenario1_traffic} summarizes background HTTP traffic performance across all twenty benchmark runs (2,776 requests decentralized, 2,800 requests centralized). 

Under decentralized active mitigation, aggregate mesh latency averaged $6.74\pm 5.94$\,ms with a 95th-percentile latency of $14.29\pm 4.28$\,ms. Exactly one request error was recorded out of 2,776 transmissions (on \texttt{ratings} during Trial~3, yielding a 99.86\% success rate for \texttt{ratings} and 99.96\% overall). Inspection of Envoy proxy telemetry confirmed that this single drop was an isolated client socket timeout unrelated to defender actions; \texttt{ratings} production routing remained 100\% intact with zero quarantine rules installed. Under centralized defense, 2,800 out of 2,800 requests completed successfully (100.00\%) with a mean mesh latency of $4.69\pm 0.98$\,ms.

\subsubsection{Causal Investigation of Centralized Missed Detections}
In the recorded benchmark trials, the decentralized architecture achieved a 10/10 mitigation success rate (100\%), whereas the centralized coordinator succeeded in 7/10 trials (70\%). While this observed difference reflects real experimental measurements, with $N=10$ trials per group it does not attain statistical significance under a two-sided Fisher's exact test ($p \approx 0.2105$). Thus, this outcome should be interpreted as an observed performance characteristic under our specific polling configuration rather than a definitive proof of inherent centralized unreliability.

A detailed forensic analysis of target pod access logs for the three missed centralized runs (Trials~1, 5, and 10) revealed the exact causal mechanism:
\begin{enumerate}
    \item \textbf{Ingestion Verification:} In all three missed trials, exactly eight attack requests arrived at \texttt{details-v1-54c5dcffdf-dqkwp} and were logged by its Envoy sidecar with HTTP 200 status codes (e.g., Trial~1 request prefix \texttt{b8eeee3b} between 15:30:38.156Z and 15:30:39.930Z).
    \item \textbf{Sequential Scrape Serialization:} The centralized coordinator sequentially queries Kubernetes log APIs across 20 monitored mesh pods. With an average roundtrip latency of $\approx 250$\,ms per pod, completing a full scrape sweep across all services requires $7.5$ to $9.5$\,seconds.
    \item \textbf{Freshness Filter Invalidation (Evidence-Supported Likely Mechanism):} The attack burst duration was approximately 1.8\,seconds. In Trials~1, 5, and 10, the burst arrived immediately after the coordinator had polled \texttt{details}. By the time the coordinator finished querying the other 19 pods and returned to \texttt{details}, the elapsed time since the burst exceeded the strict 8.0-second time-to-live threshold (\texttt{L7\_PREDICTION\_TTL}). While explicit record-discard statements were emitted at the debug level and thus filtered by the controller's default \texttt{INFO} logging configuration, this 8.0-second TTL invalidation is the evidence-supported likely mechanism, corroborated by target pod access log timestamps, scrape loop timing, and subsequent controller logs recording expired stale predictions without refresh (\texttt{[L7-EXPIRY] Expired stale prediction for details-v1 (age=9.7s > 8.0s)}).
\end{enumerate}

In contrast, each decentralized defender monitors exclusively its own service replicas (exactly 3 pods for \texttt{details}), completing its raw Kubernetes API log collection sweep in approximately 750\,ms ($3 \times 250$\,ms, with local in-memory parsing requiring $<50$\,ms). Telemetry is evaluated within $1.2$ to $2.5$\,seconds of arrival ($\ll 8.0$\,s TTL), ensuring complete burst capture. Under the evaluated configuration, decentralized control achieved lower conditional mitigation latency and detected more short-duration attack bursts than centralized control, although the observed difference in mitigation success rates was not statistically significant at the 5\% level.
